\chapter*{读法}

读这本书时，不要只看最终代码。每道题至少走四遍：第一遍只读题并写暴力思路；第二遍找慢点和可优化性质；第三遍写 C++ 优化代码；第四遍复盘边界、复杂度和面试表达。

每天训练时，把题目记录成固定格式：题型、输入规模、暴力枚举对象、优化来源、不变量、使用的数据结构、复杂度、错因和一周后是否能独立重写。只写“没想到”没有意义，必须写成具体原因：没有发现单调性、哈希表插入顺序错、窗口收缩条件错、二分边界没有排除 \code{mid}、\code{int} 溢出、递归深度过大、容器选择导致复杂度不对。

本书目录只显示部和章，小节不进入目录。章内标题用于串联正文，不把内容切碎成大量编号。你应该把每章当成一条推理链来读，而不是把每个小标题当成孤立知识点。
